\section{Equalizers and commuting pullbacks}\label{sec:setup}

The automorphism $H$ has inverse
\[
 H^{-1}(x,y)=\bigl(a^{-1}(x^q+g(x)-y),x\bigr).
\]
All its coefficients are fixed by $q$-Frobenius, so $H\Phi_q=\Phi_q H$.
Consequently
\begin{equation}\label{eq:explicit-S}
 S(x,y)=\bigl(a^{-1}(x^{q^2}+g(x^q)-y^q),x^q\bigr),
 \qquad S^n=H^{-n}\Phi_q^n.
\end{equation}
In particular $S$ is a finite radicial map: coordinatewise Frobenius
is finite and radicial, and composing with an automorphism preserves
these properties. The universal-homeomorphism property of Frobenius
is standard; see~\cite[Tag 0CC8]{StacksFrobenius}. Radiciality does
not mean that a fixed-point scheme must carry the degree of a
geometric fiber as its point count.

Composing the equalizer equations with $H^n$ identifies the fixed
scheme of $S^n$ with
\begin{equation}\label{eq:equalizer}
 \mathcal E_n=\{H^n(P)=\Phi_q^n(P)\}.
\end{equation}
Equivalently, the two pairs of coordinate differences generate the
same equalizer ideal after this automorphism: differences of
polynomial substitutions lie in the ideal generated by the original
coordinate differences, and $H^{-n}$ gives the reverse inclusion.

Let $R=\F_q[x,y]$, and introduce the $\F_q$-linear operators
\begin{equation}\label{eq:operators}
 T=H^*,\qquad U=\Phi_q^*,\qquad \delta=T-U.
\end{equation}
Here pullback means substitution; $T$ and $U$ are ring homomorphisms,
whereas their difference $\delta$ is only an additive linear
operator. In particular $\delta^j$ is an operator iterate, not an
iterate of a point map. The commutation above gives $TU=UT$.
Moreover $U(P)=P^q$ for $P\in R$.

Write $u_0=g(y)-ax$. On the two coordinates,
\begin{equation}\label{eq:delta-coordinates}
 \delta x=y-x^q,\qquad \delta y=u_0,
 \qquad Ty=y^q+u_0,
\end{equation}
and for every $P\in R$,
\begin{equation}\label{eq:ordinary-degree-bound}
 \deg\delta P\le q\deg P.
\end{equation}
We use total degree and put $\deg0=-\infty$. A polynomial has a
\emph{unique top term} $d_0y^D$ if its difference from $d_0y^D$ has
total degree strictly less than $D$, where $d_0\ne0$.

The degree calculation needed for all periods is the following more
precise version of \eqref{eq:degree-formulas}:
\begin{equation}\label{eq:leading-claim}
 \delta^j y=c_jy^{D_j}+\text{terms of degree }<D_j,
 \qquad
 c_j=
 \begin{cases}
 b(c\bar\ell)^{j-1},&\text{high support},\\
 2^{j-1}b^{(p^j-1)/(p-1)},&\text{low support}.
 \end{cases}
\end{equation}
In this expression $\bar\ell$ denotes the residue of $\ell$ in
$\F_p^*$. All $c_j$ are nonzero. We first explain why this one
sequence of leading terms determines every ordinary fixed count.
